% !TEX TS-program = pdflatex
% !TeX program = pdflatex
% !TEX encoding = UTF-8
% !TeX spellcheck = en_US

\documentclass[11pt, a4paper]{article}
%\usepackage{fullpage}
\usepackage[left=1cm,right=1cm,top=1cm,bottom=2cm]{geometry}
\usepackage[fleqn]{amsmath}
\usepackage{amssymb}
%\usepackage{indentfirst}
\usepackage[T1]{fontenc}
\usepackage[utf8]{inputenc}
\usepackage[french,english]{babel}
\usepackage{txfonts} 
\usepackage[]{graphicx}
\usepackage{multirow}
\usepackage{hyperref}
\usepackage{parskip}
\usepackage{multicol}
\usepackage{wrapfig}
\usepackage{multicol}

\usepackage{turnstile}%Induction symbole

\usepackage{tikz}
\usetikzlibrary{arrows, automata}
\usetikzlibrary{decorations.pathmorphing}

\renewcommand{\baselinestretch}{1}

\setlength{\parindent}{0pt}


\begin{document}
	
	%\pagestyle{empty} 
	\begin{tabular}{ll}
		\multirow{3}{*}{\includegraphics[width=1.5cm]{../img/esi.mlsid.pdf}} 
		& \'Ecole national Supérieure d'Informatique, Algiers, Algeria\\
		& 2CSSID (2024/2025)\\
		& Machine Learning (ML)
	\end{tabular}\\[.25cm]
	\noindent\rule{\textwidth}{1pt}\\[-0.5cm]
	\begin{center}
		{\LARGE \textbf{Workshop 03: Neural Networks' tools (Tensorflow)}}
		\begin{flushright}
			Mr. Abdelkrime ARIES
		\end{flushright}
	\end{center}\vspace{-.25cm}
	\noindent\rule{\textwidth}{1pt}
	
	\begin{center}
		\begin{tabular}{|p{.8\textwidth}|}
			\hline
			Get familiar with a tool to create MLPs and train them.  \\
			Learn how to create custom MLPs using this tool. \\
			Understand how this tool works by programming low level MLPs \\
			\hline
		\end{tabular}
	\end{center}

%\section*{Useful links}
%
%\begin{itemize}
%	\item \textbf{Dataset}: {\scriptsize \url{https://archive.ics.uci.edu/dataset/601/ai4i+2020+predictive+maintenance+dataset}} 
%	\item \textbf{Demo}: {\scriptsize \url{https://github.com/projeduc/ESI_ML/blob/main/TP/SD_2022-2023/TP01/TP01_pretraitement_sujet.ipynb}}
%	\item \textbf{MatplotLib}: {\scriptsize \url{https://matplotlib.org/stable/gallery/index.html}}
%	\item \textbf{imbalanced-learn}: {\scriptsize \url{https://imbalanced-learn.org/stable/user_guide.html#user-guide}}
%	\item \textbf{CRISP-DM (web)}: {\scriptsize \url{https://www.ibm.com/docs/en/spss-modeler/18.5.0?topic=dm-crisp-help-overview}}
%	\item  \textbf{CRISP-DM (pdf)}: {\scriptsize \url{https://www.ibm.com/docs/en/SS3RA7_18.5.0/pdf/ModelerCRISPDM.pdf}}
%\end{itemize}


\begin{multicols}{2}
	\section{Data Preparation}
	
	We will use the \textbf{satellite} dataset:
	\begin{itemize} 
		\item Import the training dataset. 
		\item Separate the input data and normalize it by dividing by \textbf{255}. 
		\item Separate the output data and encode it in OneHot format using \textbf{LabelBinarizer} from \textbf{scikit-learn}.
		\item Transform \textbf{X\_train} and \textbf{Y\_train} to \textbf{tf.constant} with \textbf{dtype=tf.float32}.
		\item Redo the same with test dataset, except do not transform \textbf{Y\_test} to \textbf{tf.constant}.
	\end{itemize}
	
	\section{Keras}
	
	In this section, we will create a project to classify satellite images.
	We will try to use high level APIs to create an MLP.
	
	\subsection{Sequential model}
	
	Create a sequential model (\textbf{nn1}). Then, add:
	\begin{itemize} 
		\item an input layer having a shape (number of features of the past dataset, );
		\item a dense layer having 10 neurons and an activation function 'relu';
		\item a dense layer having 10 neurons and an activation function 'relu';
		\item and a dense layer having number-of-classes neurons and an activation function 'softmax'.
	\end{itemize}
	Use its method \textbf{summary} to print its structure.
	
	\subsection{Model training}
	
	\begin{itemize}  
		\item Compile \textbf{nn1} by calling the method \textbf{compile}:
		\begin{itemize} 
			\item Use cross entropy as loss function.
			\item Define Adam as an optimization function.
		\end{itemize}
		\item Train \textbf{nn1} using its method \textbf{fit}.
		\begin{itemize} 
			\item Use \textbf{X\_train} and \textbf{Y\_train} as its obligatory arguments.
			\item Use a number of epochs of your choice.
		\end{itemize}
	\end{itemize}
	
	\subsection{Model testing}
	\begin{itemize}  
		\item Generate the estimations of \textbf{X\_test} using \textbf{nn1}.
		\item Transform these estimations into classes using the method \textbf{inverse\_transform} of our \textbf{LabelBinarizer}.
		\item Use these to print a classification report.
	\end{itemize}
	
	\section{High level with a custom class}
	
	We want to create a custom model which adds layers one after the other.
	It verifies if the layers are compatibles.
	Then, when finished, we cannot add any more layers (we lock it).
	
	\subsection{Custom Layer}
	
	\begin{itemize}  
		\item Create a class called \textbf{MyLayer} which inherits from \textbf{keras.layers.Dense}.
		\item Its obligatory arguments are: the number of inputs and the number of outputs.
		\item Its optional arguments are: a boolean for bias (True by default) and a string of activation function to choose from 'relu', 'sigmoid' and 'linear' ('linear' by default).
		\item Use \textbf{assert} to verify that both the number of inputs and outputs are more than 0.
		\item In this case the number of outputs is the number of units (neurons).
		To specify the number of inputs, call the method \textbf{build((number of inputs, ))}.
		\item You have some unitary tests, adapt them to your model to test it.
	\end{itemize}
	
	\subsection{Custom Net}
	
	\begin{itemize}  
		\item Create a class called \textbf{MyMLP} which inherits from \textbf{keras.Model}.
		\item It has no arguments; but, it contains 2 attributes: a lock initialized to False and a list to store our layers. 
		\item Add a method \textbf{add\_layer} which has one argument: a layer. 
		\begin{itemize} 
			\item It raises an exception when the object is locked.
			\item It raises an exception when the previous layer's output is not as the current layer's input.
			\item In case this is the first layer to add, no need to check.
			\item It adds the layer when all is right.
			\item Then, it returns \textbf{self}; this allows us to call this method in one line 
		\end{itemize}
		\item Add a method \textbf{compile} which has these optional arguments: number of inputs of the output layer (=1), number of outputs of the output layer (=1), bias of the output layer (=True), a boolean indicating if it is a multi-class classification (=False), and learning rate (=1.).
		\begin{itemize} 
			\item If there are past layers, it modifies the number of inputs of the output layer to be the output of the last layer.
			\item The output layer will have a sigmoid activation function
			\item If it is multi-class and the number of outputs is more than 1, this function must be softmax.
			\item The loss function must be BCE.
			\item If it is multi-class and the number of outputs is more than 1, The loss function must be CE.
			\item The optimizer must be Adam, taking this model's parameters and the bias.
			\item Lock the model.
		\end{itemize}
		\item Add the \textbf{forward} method taking \textbf{X} as argument.
		\item Override \textbf{\_\_call\_\_} which returns the backward pass's result.
	\end{itemize}
	
	\subsection{Model training} 
	
	\begin{itemize}  
		\item Create a model (\textbf{nn2}) composed of:
		\begin{itemize} 
			\item a hidden layer having a shape (number of features of the past dataset, 10) and ReLU activation function;
			\item a hidden layer having a shape (10, 10) and a ReLU activation function;
			\item an output layer having a shape (10, number of classes) for multi-class classification.
		\end{itemize}
		\item Call \textbf{summary} to print its structure.
		\item Fit it on \textbf{satellite} train dataset.
	\end{itemize}
	
	\subsection{Model testing} 
	
	\begin{itemize}  
		\item Generate the estimations of \textbf{X\_test} using \textbf{nn2}.
		\item Transform these estimations into classes using the method \textbf{inverse\_transform} of our \textbf{LabelBinarizer}.
		\item Use these to print a classification report.
	\end{itemize}
	
	\section{Low Level}
	
	\subsection{Activation functions}
	
	\begin{itemize}  
		\item Create functions \textbf{simple\_sigmoid}, \textbf{simple\_ReLU} and \textbf{simple\_softmax} which takes \textbf{X} as argument.
		\item You have to use \textbf{tf} functions.
	\end{itemize}
	
	\subsection{Loss functions}
	
	\begin{itemize}  
		\item Create classes \textbf{SimpleBCE} and \textbf{SimpleCE} which inherit from \textbf{keras.Loss}.
		\item No constructor is required. 
		\item Add the \textbf{call} method which takes \textbf{X, Y} as arguments.
	\end{itemize}
	
	\subsection{Optimization functions}
	
	\begin{itemize}  
		\item Create a class \textbf{SimpleGD} which inherits from \textbf{keras.Optimizer}.
		\item The constructor takes \textbf{lr} as argument. 
		\item It passes it to the super class as \textbf{super().\_\_init\_\_(learning\_rate=lr)}
		\item Override the \textbf{apply\_gradients} method which takes one argument: \textbf{grads\_and\_vars}. 
		\begin{itemize}  
			\item The argument is a list of tuples (grads, vars).
			\item Update the parameters (vars) using their method \textbf{assign\_sub} which is the original value minus the value in the arguments.
		\end{itemize}
	\end{itemize}
	
	\subsection{Custom Layer}
	
	\begin{itemize}  
		\item Copy \textbf{MyLayer} and rename it \textbf{SimpleLayer}, but it must inherit from \textbf{object}.
		\item Create an attribute \textbf{trainable\_weights} which is a list to store the variables.
		\item Create an attribute \textbf{W} of type \textbf{tf.Variable} initialized to zeros.
		Add it to \textbf{trainable\_weights}.
		\item Create an attribute \textbf{W} initialized to zeros. 
		It can be only transformed to \textbf{tf.Variable} and stored in \textbf{trainable\_weights} if bias=True.
		In this case, it will be subscribed into the module's parameters; thus can be recovered using the method \textbf{parameters}.
		\item Replace all activation functions by our simple versions.
		\item Add a method \textbf{randomize} which randomizes all trainable parameters using a normal law (\textbf{tf.random.normal}) having a mean of 0 and a standard deviation of 1.
		Try to update them using \textbf{assign} method.
		\item Add the \textbf{forward} method.
		\item Override \textbf{\_\_call\_\_} to return the result of the past method.
	\end{itemize}
	
	\subsection{Custom Net}
	
	\begin{itemize}  
		\item Copy \textbf{MyMLP} and rename it into \textbf{SimpleMLP}, but it must inherit from \textbf{object}.
		\item Create an attribute \textbf{trainable\_weights} which is a list to store the variables.
		\item Modify it to handle our new structures.
		\item Add a method \textbf{randomize} which randomizes all layers.
		\item Add a method \textbf{compile} which has these optional arguments: number of inputs of the output layer (=1), number of outputs of the output layer (=1), bias of the output layer (=True), a boolean indicating if it is a multi-class classification (=False), and learning rate (=1.).
		\begin{itemize} 
			\item If there are past layers, it modifies the number of inputs of the output layer to be the output of the last layer.
			\item The output layer will have a sigmoid activation function
			\item If it is multi-class and the number of outputs is more than 1, this function must be softmax.
			\item The loss function must be BCE.
			\item If it is multi-class and the number of outputs is more than 1, The loss function must be CE.
			\item The optimizer must be Adam, taking this model's parameters and the bias.
			\item Lock the model.
			\item Update the \textbf{trainable\_weights}.
		\end{itemize}
		\item Add the \textbf{forward} method taking \textbf{X} as argument.
		\item Add the \textbf{backward} method taking \textbf{X, Y} as arguments:
		\begin{itemize} 
			\item In the context of a \textbf{tf.GradientTape()} called \textbf{tape}, perform a forward pass to get the predicted classes. Calculate the error using the \textbf{loss} function stored as an attribute. 
			\item Exit the \textbf{tape} context and compute the gradients using \textbf{tape.gradient} with the error and the list of parameters as arguments. 
			\item Use the optimizer to apply the gradients (method \textbf{apply\_gradients}) which takes the gradients and the zipped list of parameters as arguments. 
			\item Return the overall error as numpy with \textbf{.numpy()}. 
		\end{itemize}
		\item Add a \textbf{fit} method taking \textbf{X, Y, epochs} as arguments.
		For each iteration, it shows the cost.
		\item Override \textbf{\_\_call\_\_} which returns the backward pass's result. 
	\end{itemize}
	
	\subsection{Model training} 
	
	\begin{itemize}  
		\item Create a model (\textbf{nn3}) composed of:
		\begin{itemize} 
			\item a hidden layer having a shape (number of features of the past dataset, 10) and ReLU activation function;
			\item a hidden layer having a shape (10, 10) and a ReLU activation function;
			\item an output layer having a shape (10, number of classes) for multi-class classification.
		\end{itemize}
		\item Randomize it.
		\item Print its parameters. The method \textbf{parameters} returns an iterator, so you have to cast it into a list.
		\item Fit it on \textbf{satellite} train dataset.
	\end{itemize}
	
	\subsection{Model testing} 
	
	\begin{itemize}  
		\item Generate the estimations of \textbf{X\_test} using \textbf{nn3}.
		\item Transform these estimations into classes using the method \textbf{inverse\_transform} of our \textbf{LabelBinarizer}.
		\item Use these to print a classification report.
	\end{itemize}

\end{multicols}

\section*{Appendix}

\begin{itemize}  
	\item pandas.read\_csv: \url{https://pandas.pydata.org/docs/reference/api/pandas.read_csv.html}
	\item sklearn.preprocessing.LabelBinarizer: \url{https://scikit-learn.org/1.5/modules/generated/sklearn.preprocessing.LabelBinarizer.html}
	\item tf.constant: \url{https://www.tensorflow.org/api_docs/python/tf/constant}
	\item tf.keras.Sequential: \url{https://www.tensorflow.org/api_docs/python/tf/keras/Sequential}
	\item tf.keras.Input: \url{https://www.tensorflow.org/api_docs/python/tf/keras/Input}
	\item tf.keras.layers.Dense: \url{https://www.tensorflow.org/api_docs/python/tf/keras/layers/Dense}
	\item tf.keras.optimizers.Adam: \url{https://www.tensorflow.org/api_docs/python/tf/keras/optimizers/Adam}
	\item tf.keras.losses.CategoricalCrossentropy: \url{https://www.tensorflow.org/api_docs/python/tf/keras/losses/CategoricalCrossentropy}
	\item sklearn.metrics.classification\_report: \url{https://scikit-learn.org/1.5/modules/generated/sklearn.metrics.classification_report.html}
	\item assert: \url{https://docs.python.org/3/reference/simple_stmts.html#assert}
	\item tf.keras.Model: \url{https://www.tensorflow.org/api_docs/python/tf/keras/Model}
	\item Exception: \url{https://docs.python.org/3/library/exceptions.html#Exception}
	\item tf.keras.losses.BinaryCrossentropy: \url{https://www.tensorflow.org/api_docs/python/tf/keras/losses/BinaryCrossentropy}
	\item tf.keras.Loss: \url{https://www.tensorflow.org/api_docs/python/tf/keras/Loss}
	\item tf.math.exp: \url{https://www.tensorflow.org/api_docs/python/tf/math/exp}
	\item tf.where: \url{https://www.tensorflow.org/api_docs/python/tf/where}
	\item tf.reshape: \url{https://www.tensorflow.org/api_docs/python/tf/reshape}
	\item tf.reduce\_sum: \url{https://www.tensorflow.org/api_docs/python/tf/math/reduce_sum}
	\item tf.reduce\_mean: \url{https://www.tensorflow.org/api_docs/python/tf/math/reduce_mean}
	\item tf.math.log: \url{https://www.tensorflow.org/api_docs/python/tf/math/log}
	\item tf.keras.Optimizer: \url{https://www.tensorflow.org/api_docs/python/tf/keras/Optimizer}
	\item tf.Variable: \url{https://www.tensorflow.org/api_docs/python/tf/Variable}
	\item tf.zeros: \url{https://www.tensorflow.org/api_docs/python/tf/zeros}
	\item tf.random.normal: \url{https://www.tensorflow.org/api_docs/python/tf/random/normal}
	\item tf.matmul: \url{https://www.tensorflow.org/api_docs/python/tf/linalg/matmul}
	\item tf.GradientTape: \url{https://www.tensorflow.org/api_docs/python/tf/GradientTape}
	
%	\item tf.keras.Layer: \url{https://www.tensorflow.org/api_docs/python/tf/keras/Layer}
%	\item tf.keras.activations.sigmoid: \url{https://www.tensorflow.org/api_docs/python/tf/keras/activations/sigmoid}
\end{itemize}


\end{document}
